
In this paper, I contribute to two broad stands of economic literature. My model contributes to existing theoretical work on models of occupational choice and firm size. My empirical section contributes to a vast literature on the effects of migration in the United States.

\subsection{Theoretical Models}

    I build upon a literature of Span-of-Control Models since the seminal work of \cite{LucasModel}. In this class of models, agents are endowed with heterogeneous managerial talent. Based on this talent amount, agents sort into entrepreneurs/managers and workers. The endogenous firm-size distribution is pinned down by the underlying talent distribution of the economy in a static equilibrium. Lucas's model provides a clean and tractable framework to understand why firms differ in size and who chooses to become an entrepreneur. \cite{Azoulay} extend this model by allowing for heterogeneous underlying talent distributions for Immigrants versus native Americans. The authors find evidence that immigrants are positively selected on managerial talent. More precisely, the authors find a higher right tail of talent for immigrant entrepreneurs. In other words, looking at the largest firms in the US, immigrant founders are disproportionately overrepresented. My model is focused on identifying the mass of talent below the endogenous entry-exit threshold. This region informs us about the marginal/unsuccessful entrepreneurs who drive most of the churn in the economy. My findings are consistent with positive selection of immigrants, wherein natives have more mass below a an entry/exit threshold than immigrants.
    
    I extend this line of work by incorporating a few novel ideas. I account for credit constraints on new businesses. These constraints limit how much labor an entrepreneur can hire and the subsequent firm size. This accounts for the fact that super-talented individuals may not realize their true potential due to limits on how much money they have access to. I also include talent uncertainty before entry into business. Individuals are not aware of their true managerial talent before at least one-period of production. This model builds upon previous work on learning about one's own ability, such as \cite{Jovanovic}. My model adapts this framework to study the specific question of studying immigrant versus native entrepreneurs. I used a cross-sectional sample of U.S. businesses to infer characteristics of the underlying talent distribution of natives versus immigrants.
    
\subsection{Empirical Work}

    A large swath of literature in economics examines the effects of immigration. Recent empirical work generally finds positive effects of immigration across various aspects, such as Productivity \citep{EffectofImmigrationProductivity}, patenting and wages \citep{ImmigrationInnovationGrowth}, and  Foreign Direct Investment \citep{MigrantsAncestorsandInvestments}. \cite{ImmigrationAndBusinessDynamics} shows that US immigration since 2000 causes an increase in employment and the number of businesses at the commuting zone level. He shows that these effects are mostly driven by a reduction in exit rate. In my paper, I expand towards looking at different measures of business dynamism, such as firm size. I conduct my analysis at the more granular level of US counties and look at the longer time period of 1975 to 2010.  In this study, I utilize the identification strategy of \cite{ImmigrationInnovationGrowth} to deal with the endogeneity concerns of immigration choice. The authors build upon the canonical \cite{CardShiftShare} shift-share instruments. However, instead of instrumenting migration with realized ancestry, they instrument using predicted ancestry. This provides a more robust instrument to study immigration. This study provides a new dimension in studying the effects of immigration by using a relatively new identification technique and a more granular setting of US counties than previous work. It also contributes to the widely discussed topic of the slowdown in U.S. business dynamism by examining the aspect of immigration \citep{TenFactsOnDecliningBusinessDynamism, AreIdeasGettingHarderToFind, bessen2016, andrews2016, PopulationGrowthFirmDynamics, FromPopulationGrowthtoFirmDemographics}. 
    
    
    